%% ======================================================================
%% chap7/results_main.tex
%%
%% Chapter 7: Results
%%
%% Reports the measured results from all evaluation activities. Technical
%% results cover throughput, latency, block weight utilisation, storage
%% growth, reliability, the Snowbridge end-to-end bridge demonstration,
%% and the comparative analysis against a centralized Express.js + SQLite
%% baseline. Economic results present the Monte Carlo creator revenue
%% simulation, the audience-size crossover finding, and the fee analysis.
%% Security and user-centric results cover the security audit, test
%% coverage, decentralization (HHI) assessment, and static analysis.
%% Interpretation is deferred to Chapter 8.
%%
%% Sections:
%%   7.1 Technical Results (7 subsections, Tables 7.1-7.7)
%%   7.2 Economic Results (3 subsections, Tables 7.8-7.12)
%%   7.3 Security and User-Centric Results (4 subsections, Tables 7.13-7.14)
%%   7.4 Summary
%%
%% Tables (14 total): all use longtable with L{} and R{} fixed-width
%% columns to stay within the text block on both page sides.
%%
%% External labels referenced:
%%   chap:discussion (Ch 8), sec:limitations (8.5),
%%   sec:test-scenarios (6.4),
%%   app:challenges (Appendix A), app:security (Appendix D)
%%
%% Citations: none (empirical results chapter; sources are our own
%% measurements and the Monte Carlo simulation).
%% ======================================================================

\chapter{Results}\label{chap:results}

This chapter presents the measured results. Section~\ref{sec:technical-results} covers technical results, Section~\ref{sec:economic-results} covers economic results, and Section~\ref{sec:security-user-results} covers security and user-centric results. We defer interpretation to Chapter~\ref{chap:discussion}.

\section{Technical Results}\label{sec:technical-results}

The technical results address SQ1 (Section~\ref{subsec:latency}), SQ2 (Section~\ref{subsec:comparative-analysis}), and SQ3 (Sections~\ref{subsec:throughput}--\ref{subsec:storage}).

\subsection{Throughput}\label{subsec:throughput}

Table~\ref{tab:throughput} reports 100 concurrent extrinsics per type against a single collator (6-second blocks). All four non-minting extrinsics achieved 100/100 inclusion in one block ($\geq$100 extrinsics per block, $\approx$16.6~TPS). No failure arose from block weight, space, or transaction-pool capacity.

All failures were dispatch errors from the fifty-child bound (Section~\ref{subsec:nft-tokenization}): \texttt{MaxChildrenReached} directly, or cascading \texttt{SubscriptionNotFound}/\texttt{ViewPackNotFound} for dependent extrinsics. The benchmark characterizes storage-design capacity, not parachain throughput. Sustained throughput is in Section~\ref{subsec:stress-test}.

Per-extrinsic TPS is not reported: elapsed time reflects how many blocks the harness waited before observing inclusion (1--4 blocks), not execution cost. Relative costs are in the FRAME weights (Appendix~\ref{app:code}; hand-estimated, not benchmarked).

\begin{longtable}{@{}L{0.26\linewidth} R{0.10\linewidth} R{0.10\linewidth} R{0.09\linewidth} L{0.35\linewidth}@{}}
\caption{Single-Chain Inclusion Capacity at Batch Size 100}\label{tab:throughput}\\
\toprule
\textbf{Extrinsic} & \textbf{Succ.} & \textbf{Failed} & \textbf{Blocks} & \textbf{Failure cause} \\
\midrule
\endfirsthead
\multicolumn{5}{l}{\emph{Table~\ref{tab:throughput} continued from previous page}}\\
\toprule
\textbf{Extrinsic} & \textbf{Succ.} & \textbf{Failed} & \textbf{Blocks} & \textbf{Failure cause} \\
\midrule
\endhead
\midrule
\multicolumn{5}{r}{\emph{Continued on next page}}\\
\endfoot
\bottomrule
\endlastfoot

\texttt{register\_content}      & 100 &   0 & 1 & --- \\
\texttt{check\_access}          & 100 &   0 & 1 & --- \\
\texttt{set\_royalty\_splits}   & 100 &   0 & 1 & --- \\
\texttt{query\_rights\_metadata} & 100 &  0 & 1 & --- \\
\addlinespace
\texttt{subscribe}             &  50 &  50 & 1 & \texttt{MaxChildrenReached} \\
\texttt{purchase\_ownership}    &  50 &  50 & 1 & \texttt{MaxChildrenReached} \\
\texttt{purchase\_views}        &   0 & 100 & 0 & \texttt{MaxChildrenReached} \\
\addlinespace
\texttt{renew\_subscription}    &  50 &  50 & 1 & \texttt{SubscriptionNotFound}$^{\dagger}$ \\
\texttt{enable\_auto\_renew}    &  50 &  50 & 1 & \texttt{SubscriptionNotFound}$^{\dagger}$ \\
\texttt{disable\_auto\_renew}   &  50 &  50 & 1 & \texttt{SubscriptionNotFound}$^{\dagger}$ \\
\texttt{consume\_view}          &  22 &  78 & 1 & \texttt{ViewPackNotFound}$^{\dagger}$ \\
\end{longtable}

{\footnotesize $\dagger$~Cascading consequence of the fifty-child bound: the prerequisite subscriptions and view packs could not be created, so the dependent extrinsics had no state to act on.}

\subsection{Sustained Throughput (Stress Test)}\label{subsec:stress-test}

Under sustained load (200 accounts, 180~s), the system achieved 3,875 successful extrinsics at a \textbf{96.9\% success rate} over the first 142~s, yielding \textbf{27.3~TPS} (174.6 transactions per block). The overall failure count is a harness artifact: nonce exhaustion caused the benchmark client to spin thousands of submissions in six seconds at 142--148~s; this was a client limit, not chain capacity.

\subsection{Latency}\label{subsec:latency}

Single-chain inclusion latency was consistently \textbf{$\sim$6,000 ms} (one parachain block; s.d.\ $<$15~ms), confirming atomic same-block inclusion. XCM cross-chain latencies are in Table~\ref{tab:xcm-latency} ($n=3$ per operation).

\begin{longtable}{@{}L{0.32\linewidth} R{0.18\linewidth} R{0.18\linewidth} R{0.24\linewidth}@{}}
\caption{XCM Destination-Chain Inclusion Delay by Operation ($n=3$ per operation)}\label{tab:xcm-latency}\\
\toprule
\textbf{Operation} & \textbf{Mean Block Delta} & \textbf{Standard Deviation} & \textbf{15-block scan window} \\
\midrule
\endfirsthead

\multicolumn{4}{l}{\emph{Table~\ref{tab:xcm-latency} continued from previous page}}\\
\toprule
\textbf{Operation} & \textbf{Mean Block Delta} & \textbf{Standard Deviation} & \textbf{15-block scan window} \\
\midrule
\endhead

\midrule
\multicolumn{4}{r}{\emph{Continued on next page}}\\
\endfoot

\bottomrule
\endlastfoot

\texttt{xcm\_subscribe} & 3.0 blocks & $\pm$2.0 blocks & $\sim$99.6 seconds \\
\texttt{xcm\_purchase\_views} & 5.3 blocks & $\pm$1.5 blocks & $\sim$100 seconds \\
\texttt{xcm\_purchase\_ownership} & 5.3 blocks & $\pm$1.5 blocks & $\sim$100 seconds \\
\end{longtable}

The destination event was observed within 15 blocks in all nine runs (median delta 5, $\sim$30--35~s at 6.67~s/block); the $\sim$100~s wall-clock measures producing 15 destination blocks. Async Backing and Agile Coretime \citep{parityUpgrade2025} would reduce these delays.

\subsection{Block Weight Utilization}\label{subsec:block-weight}

\begin{longtable}{@{}R{0.16\linewidth} R{0.25\linewidth} R{0.25\linewidth} R{0.26\linewidth}@{}}
\caption{Block Weight Utilization at Increasing Batch Sizes}\label{tab:block-weight}\\
\toprule
\textbf{Batch Size} & \textbf{\texttt{ref\_time} Consumed} & \textbf{\texttt{proof\_size} Consumed} & \textbf{Block Construction Time} \\
\midrule
\endfirsthead

\multicolumn{4}{l}{\emph{Table~\ref{tab:block-weight} continued from previous page}}\\
\toprule
\textbf{Batch Size} & \textbf{\texttt{ref\_time} Consumed} & \textbf{\texttt{proof\_size} Consumed} & \textbf{Block Construction Time} \\
\midrule
\endhead

\midrule
\multicolumn{4}{r}{\emph{Continued on next page}}\\
\endfoot

\bottomrule
\endlastfoot

50 & 2.2\% & 0.4\% & $\sim$10 ms \\
100 & 4.3\% & 0.7\% & $\sim$21 ms \\
150 & 6.5\% & 1.1\% & $\sim$31 ms \\
200 & 8.6\% & 1.5\% & $\sim$42 ms \\
250 & 10.7\% & 1.8\% & $\sim$52 ms \\
300 & 12.9\% & 2.2\% & $\sim$62 ms \\
\end{longtable}

At 300 transactions, block weight accounts for 12.9\% of \texttt{ref\_time} and 2.2\% of \texttt{proof\_size}, with a construction time of approximately 62~ms ($\sim$1\% of the 6-second block). PoV weight is zero in every hand-estimated weight (Appendix~\ref{app:code}); this figure reflects reference time only.

\subsection{Storage Growth}\label{subsec:storage}

\begin{longtable}{@{}L{0.34\linewidth} L{0.28\linewidth} R{0.30\linewidth}@{}}
\caption{Per-Item Storage Costs (SCALE-encoded)}\label{tab:storage-costs}\\
\toprule
\textbf{Item Type} & \textbf{Storage Map} & \textbf{Bytes per Item} \\
\midrule
\endfirsthead

\multicolumn{3}{l}{\emph{Table~\ref{tab:storage-costs} continued from previous page}}\\
\toprule
\textbf{Item Type} & \textbf{Storage Map} & \textbf{Bytes per Item} \\
\midrule
\endhead

\midrule
\multicolumn{3}{r}{\emph{Continued on next page}}\\
\endfoot

\bottomrule
\endlastfoot

Content registration & \texttt{Contents} & 191 \\
Active subscription & \texttt{Subscriptions} & 112 \\
View pack & \texttt{ViewPacks} & 111 \\
Ownership record & \texttt{Ownership} & 36 \\
Royalty splits (per collaborator) & \texttt{RoyaltySplits} & 34 \\
Child NFT nesting record & \texttt{Children} & 8 (per child) \\
\end{longtable}

Growth was perfectly linear (the 50th item costs the same as the first), confirming \textbf{$O(1)$} \texttt{StorageMap} complexity. At one million content items and ten million subscribers, projected state is approximately 1.3~GB, within the 500-byte-per-record Phase~2 target.

\subsection{Reliability}\label{subsec:reliability}

\begin{longtable}{@{}L{0.16\linewidth} L{0.42\linewidth} L{0.35\linewidth}@{}}
\caption{Reliability Test Results}\label{tab:reliability}\\
\toprule
\textbf{Phase} & \textbf{Metric} & \textbf{Value} \\
\midrule
\endfirsthead

\multicolumn{3}{l}{\emph{Table~\ref{tab:reliability} continued from previous page}}\\
\toprule
\textbf{Phase} & \textbf{Metric} & \textbf{Value} \\
\midrule
\endhead

\midrule
\multicolumn{3}{r}{\emph{Continued on next page}}\\
\endfoot

\bottomrule
\endlastfoot

Baseline & Blocks produced & 18 of 20 expected \\
Baseline & Uptime & 90\% \\
Baseline & Mean block time & 6.67 seconds \\
Failure & Time to RPC ready (post-restart) & $\sim$15 seconds \\
Failure & Time to first new block (post-restart) & $\sim$20 seconds \\
Recovery & State integrity & 100\% preserved \\
Recovery & Blocks produced & 18 of 20 expected \\
Recovery & Uptime & 90\% (matches baseline) \\
Recovery & Mean block time & 6.67 seconds (matches baseline) \\
\end{longtable}

A 90\% block production rate reflects typical single-collator Zombienet behavior (Finding~F); the collator did not go down. MTTR was approximately 15~s for RPC and 20~s for the first new block, with 100\% state preservation. Recovery metrics matched the baseline exactly.

\subsection{Snowbridge End-to-End Bridge}\label{subsec:snowbridge-bridge}

We bridged two 1~ETH transactions from the local Ethereum network to AssetHub via Snowbridge (four parachains, Geth, Lodestar, 16 Gateway contracts, two relayers). Both succeeded; the destination account held 22~ETH on AssetHub (20 pre-minted, 2 bridged). Table~\ref{tab:snowbridge-path} provides the message pathway.

\begin{longtable}{@{}L{0.04\linewidth} L{0.20\linewidth} L{0.70\linewidth}@{}}
\caption{Snowbridge End-to-End Message Path (ETH Transfers to AssetHub; CCRMS Not Exercised)}\label{tab:snowbridge-path}\\
\toprule
\textbf{Step} & \textbf{Component} & \textbf{Action} \\
\midrule
\endfirsthead

\multicolumn{3}{l}{\emph{Table~\ref{tab:snowbridge-path} continued from previous page}}\\
\toprule
\textbf{Step} & \textbf{Component} & \textbf{Action} \\
\midrule
\endhead

\midrule
\multicolumn{3}{r}{\emph{Continued on next page}}\\
\endfoot

\bottomrule
\endlastfoot

1 & Geth & Transaction \texttt{Gateway.sendToken(1 ETH, Alice)} included in block 5157 \\
2 & Lodestar & Block finalized at epoch $\geq$ 162 \\
3 & Beacon Relay & Generated inclusion proof, submitted beacon header update \\
4 & Bridge Hub & \texttt{ethereumBeaconClient.submit(...)} verified BLS signatures, advanced \texttt{latestFinalizedBlockRoot} \\
5 & Execution Relay & Generated receipt + execution + ancestry proof, submitted message bundle \\
6 & Bridge Hub & \texttt{ethereumInboundQueue.submit(...)} verified the proof against the verified beacon state root \\
7 & Bridge Hub $\to$ AssetHub & XCM message dispatched via HRMP \\
8 & AssetHub & \texttt{foreignAssets.mint(Ether, Alice, 1 ETH)} succeeded \\
\end{longtable}

CCRMS was not on the path; the result confirms the local pipeline and the CCRMS runtime's location-converter configuration. The final hop to a CCRMS operation was not tested end-to-end (Section~\ref{sec:limitations}); integration challenges are in Appendix~\ref{app:challenges}.

\subsection{Comparative Analysis vs.\ Centralized Baseline}\label{subsec:comparative-analysis}

We benchmarked CCRMS against Express.js with SQL.js for the same six operations, using the same methodology. Table~\ref{tab:comparative} presents the comparison.

\begin{longtable}{@{}L{0.28\linewidth} L{0.22\linewidth} L{0.22\linewidth} L{0.20\linewidth}@{}}
\caption{Comparative Analysis: CCRMS Parachain vs.\ Centralized (Express.js + SQLite)}\label{tab:comparative}\\
\toprule
\textbf{Metric} & \textbf{CCRMS Parachain} & \textbf{Centralized} & \textbf{Ratio} \\
\midrule
\endfirsthead

\multicolumn{4}{l}{\emph{Table~\ref{tab:comparative} continued from previous page}}\\
\toprule
\textbf{Metric} & \textbf{CCRMS Parachain} & \textbf{Centralized} & \textbf{Ratio} \\
\midrule
\endhead

\midrule
\multicolumn{4}{r}{\emph{Continued on next page}}\\
\endfoot

\bottomrule
\endlastfoot

Sustained TPS & 26.2 & 7,305 & $270\times$ \\
\addlinespace
Operation latency (server-side) & $\sim$6,000 ms & 0.09 ms & $67{,}000\times$ \\
\addlinespace
Stress test success rate & 22\% (RPC nonce conflicts) & 100\% & -- \\
\addlinespace
Storage per content & 191 bytes & $\sim$100--150 bytes & $1.5\times$ \\
\addlinespace
MTTR & 15--20 seconds & 2--3 seconds & $7\times$ \\
\addlinespace
Cross-chain equivalent & $\sim$100 seconds & 5--50 ms (HTTP request) & $2{,}000$--$20{,}000\times$ \\
\end{longtable}

The 22\% stress-test figure in Table~\ref{tab:comparative} is the client-visible success rate over the full saturation window, including the nonce-exhaustion burst described in Section~\ref{subsec:stress-test}. It is not the 96.9\% happy-path inclusion rate used in the abstract and Section~\ref{subsec:stress-test}, nor is it a consensus failure rate.

The centralized implementation outperforms CCRMS across all quantitative metrics ($270\times$ TPS, $67{,}000\times$ latency). CCRMS's qualitative properties (no single point of failure, cross-chain support, provable rights state, and censorship resistance) are not reflected in quantitative data but are central to the use case (gaps G1 and G7, Chapter~\ref{chap:background}). At 26.2 TPS ($\approx$2.3 million operations per day), CCRMS suffices for creator platforms where rights transactions are low-frequency events.

\section{Economic Results}\label{sec:economic-results}

The economic results address SQ5 using Monte Carlo simulation, the audience-size crossover finding, and a comparative fee analysis.

\subsection{Monte Carlo Creator Revenue Simulation}\label{subsec:monte-carlo}

The simulation models 1,000 creators across 10,000 iterations under four platform configurations, projecting one-year revenue using platform-specific fees, algorithmic discovery, and monthly churn. It uses a seeded pseudo-random number generator (xoshiro128**, seed~42), so every figure is exactly reproducible. Full parameters are in Appendix~\ref{app:monte-carlo}; sensitivity analysis in Section~\ref{subsec:sensitivity}.

\begin{figure}[htbp]
\vspace{\baselineskip}
\centering
\begin{tikzpicture}[
    flowstep/.style={draw, rectangle, rounded corners=3pt, minimum width=7.5cm,
                     minimum height=1.05cm, align=center, font=\small},
    plat/.style={draw, rectangle, rounded corners=2pt, minimum width=2.1cm,
                 minimum height=1.05cm, align=center, font=\footnotesize},
    >=Stealth
]
%% Step 1: input
\node[flowstep, fill=gray!15] (s1) at (0, 0) {
    \textbf{Creator Parameter Sampling (1,000 creators)}\\
    {\footnotesize followers $\sim$ log-normal distribution \quad churn rate \quad growth rate \quad content volume}
};
%% Step 2: iteration loop
\node[flowstep, fill=blue!12] (s2) at (0, -1.9) {
    \textbf{10,000 Revenue Iterations per Creator}\\
    {\footnotesize uniform fee draw · algorithmic discovery boost · monthly churn}
};
\draw[->, thick] (s1.south) -- (s2.north);
%% Branch point below s2
\coordinate (sp) at (0, -2.75);
\draw[thick] (s2.south) -- (sp);
%% Four platform nodes
\node[plat, fill=red!12]    (p1) at (-4.0, -3.55) {\textbf{Centralized}\\{\footnotesize 30--45\% fee}};
\node[plat, fill=orange!18] (p2) at (-1.3, -3.55) {\textbf{Web3}\\{\footnotesize 7.5--15\% fee}};
\node[plat, fill=blue!15]   (p3) at ( 1.3, -3.55) {\textbf{Bridge}\\{\footnotesize 8--15\% fee}};
\node[plat, fill=green!18]  (p4) at ( 4.0, -3.55) {\textbf{CCRMS}\\{\footnotesize 1--5\% fee}};
\draw[->, thick] (sp) -| (p1.north);
\draw[->, thick] (sp) -| (p2.north);
\draw[->, thick] (sp) -| (p3.north);
\draw[->, thick] (sp) -| (p4.north);
%% Join point above s3
\coordinate (jp) at (0, -4.6);
\draw[thick] (p1.south) |- (jp);
\draw[thick] (p2.south) |- (jp);
\draw[thick] (p3.south) |- (jp);
\draw[thick] (p4.south) |- (jp);
%% Step 3: aggregate outputs
\node[flowstep, fill=green!18] (s3) at (0, -5.5) {
    \textbf{Aggregate Distributions (per Platform)}\\
    {\footnotesize mean annual revenue · 95\% confidence intervals · per-scenario win rates}
};
\draw[->, thick] (jp) -- (s3.north);
\end{tikzpicture}
\caption{Monte Carlo simulation structure: 1,000 creators $\times$ 10,000 iterations = 10 million revenue paths per platform. Platform fee drawn uniformly each iteration; full parameters in Appendix~\ref{app:monte-carlo}.}
\label{fig:monte-carlo-flow}
\end{figure}

\begin{longtable}{@{}L{0.34\linewidth} L{0.14\linewidth} R{0.20\linewidth} L{0.24\linewidth}@{}}
\caption{Monte Carlo Creator Revenue Simulation (Mean Annual Revenue Per Creator)}\label{tab:monte-carlo}\\
\toprule
\textbf{Platform Model} & \textbf{Fee Range} & \textbf{Mean Revenue} & \textbf{vs.\ CCRMS} \\
\midrule
\endfirsthead

\multicolumn{4}{l}{\emph{Table~\ref{tab:monte-carlo} continued from previous page}}\\
\toprule
\textbf{Platform Model} & \textbf{Fee Range} & \textbf{Mean Revenue} & \textbf{vs.\ CCRMS} \\
\midrule
\endhead

\midrule
\multicolumn{4}{r}{\emph{Continued on next page}}\\
\endfoot

\bottomrule
\endlastfoot

Centralized (YouTube/Spotify-style) & 30--45\% & \$45,416 & CCRMS +16.8\% \\
\addlinespace
Web3 marketplace (OpenSea-style) & 7.5--15\% & \$48,711 & CCRMS +8.9\% \\
\addlinespace
Bridge-based cross-chain & 8--15\% & \$45,369 & CCRMS +16.9\% \\
\addlinespace
\textbf{CCRMS} & \textbf{1--5\%} & \textbf{\$53,053} & \textbf{Baseline} \\
\end{longtable}

\subsection{Audience Size Dependency}\label{subsec:audience-dependency}

CCRMS's mean advantage (Table~\ref{tab:monte-carlo}) is not uniform: it is heavily dependent on \textbf{audience size}. Table~\ref{tab:audience-tiers} breaks results down by tier.

\begin{longtable}{@{}L{0.12\linewidth} L{0.18\linewidth} R{0.20\linewidth} L{0.42\linewidth}@{}}
\caption{CCRMS Win Rate Against Centralized Platform by Audience Tier}\label{tab:audience-tiers}\\
\toprule
\textbf{Tier} & \textbf{Followers} & \textbf{Win rate vs centralized} & \textbf{Notes} \\
\midrule
\endfirsthead

\multicolumn{4}{l}{\emph{Table~\ref{tab:audience-tiers} continued from previous page}}\\
\toprule
\textbf{Tier} & \textbf{Followers} & \textbf{Win rate vs centralized} & \textbf{Notes} \\
\midrule
\endhead

\midrule
\multicolumn{4}{r}{\emph{Continued on next page}}\\
\endfoot

\bottomrule
\endlastfoot

Tiny   & $<$100            & 0.0\% & Centralized algorithmic discovery dominates \\
\addlinespace
Small  & 100--500          & 0.0\% & Fee savings insufficient to offset discovery \\
\addlinespace
Medium & 500--2,000        & 6.6\% & Transition begins \\
\addlinespace
Large  & 2,000--10,000     & 57.6\% & Crossover; the 50\% point falls near 3,600 followers \\
\addlinespace
Huge   & 10,000--100,000   & 99.0\% & CCRMS strongly preferred \\
\addlinespace
Mega   & 100,000+          & 99.9\% & CCRMS near-certain advantage \\
\end{longtable}

For creators with fewer than $\sim$3,600 followers (transition beginning near 2,000 where the win rate reaches 6.6\%), centralized discovery outweighs fee savings. For those with 10,000 or more followers, the fee differential dominates and CCRMS is preferred in over 99\% of scenarios.

These results describe the CCRMS design; the prototype's fifty-holder bound (Section~\ref{sec:limitations}) must be lifted before the mid-size and large tiers can be served.


Table~\ref{tab:breakeven} reports the per-scenario break-even rates.

\begin{longtable}{@{}L{0.55\linewidth} R{0.35\linewidth}@{}}
\caption{Per-Scenario Break-Even Rates}\label{tab:breakeven}\\
\toprule
\textbf{Comparison} & \textbf{CCRMS Per-Scenario Win Rate} \\
\midrule
\endfirsthead

\multicolumn{2}{l}{\emph{Table~\ref{tab:breakeven} continued from previous page}}\\
\toprule
\textbf{Comparison} & \textbf{CCRMS Per-Scenario Win Rate} \\
\midrule
\endhead

\midrule
\multicolumn{2}{r}{\emph{Continued on next page}}\\
\endfoot

\bottomrule
\endlastfoot

CCRMS vs Centralized & 28.0\% \\
CCRMS vs Web3 marketplace & 38.6\% \\
CCRMS vs Bridge-based cross-chain & 66.5\% \\
\end{longtable}

Per-scenario win rates are lower because CCRMS loses in small-audience scenarios; the 66.5\% win rate against bridge-based cross-chain reflects shared fee assumptions.

Table~\ref{tab:ci} reports 95\% confidence intervals for total simulated revenue.

\begin{longtable}{@{}L{0.40\linewidth} L{0.40\linewidth}@{}}
\caption{95\% Confidence Intervals for Total Simulated Revenue (1,000 Creators)}\label{tab:ci}\\
\toprule
\textbf{Platform} & \textbf{Total Revenue Range} \\
\midrule
\endfirsthead

\multicolumn{2}{l}{\emph{Table~\ref{tab:ci} continued from previous page}}\\
\toprule
\textbf{Platform} & \textbf{Total Revenue Range} \\
\midrule
\endhead

\midrule
\multicolumn{2}{r}{\emph{Continued on next page}}\\
\endfoot

\bottomrule
\endlastfoot

Centralized & \$30.9M--\$69.1M \\
Web3 marketplace & \$35.1M--\$78.4M \\
Bridge-based cross-chain & \$30.7M--\$69.4M \\
\textbf{CCRMS} & \textbf{\$36.7M--\$84.1M} \\
\end{longtable}

These intervals overlap across all platforms; the per-scenario win rates in Table~\ref{tab:audience-tiers} provide the stronger comparison.
\subsection{Comparative Fee Analysis}\label{subsec:fee-analysis}

\begin{longtable}{@{}L{0.22\linewidth} L{0.14\linewidth} L{0.56\linewidth}@{}}
\caption{Published fee ranges by platform category. Centralized literature values span 30--60\%; the simulation uses 30--45\% (Appendix~\ref{app:monte-carlo}, Table~G.2).}\label{tab:fees}\\
\toprule
\textbf{Category} & \textbf{Fee Range} & \textbf{Representative Platforms} \\
\midrule
\endfirsthead

\multicolumn{3}{l}{\emph{Table~\ref{tab:fees} continued from previous page}}\\
\toprule
\textbf{Category} & \textbf{Fee Range} & \textbf{Representative Platforms} \\
\midrule
\endhead

\midrule
\multicolumn{3}{r}{\emph{Continued on next page}}\\
\endfoot

\bottomrule
\endlastfoot

Centralized & 30--60\% & YouTube \citep{youtubePartnerProgram}, Spotify \citep{spotifyForArtists}, Apple Music \citep{appleForArtists}, Audible \citep{audibleAcx} \\
\addlinespace
Web3 marketplace & 7.5--15\% & OpenSea (2.5\% protocol fee \citep{openSeaFees}), Magic Eden (2\% \citep{magicEdenFees}) \\
\addlinespace
Bridge-based cross-chain & 8--15\% & Wormhole, LayerZero, Multichain (bridge + DEX swap fees; varies by route) \\
\addlinespace
CCRMS & 1--5\% & Polkadot transaction fees (target), conservative range \\
\end{longtable}

\subsection{Sensitivity Analysis}\label{subsec:sensitivity}

A sensitivity analysis (Appendix~\ref{app:monte-carlo}) finds the advantage insensitive to pricing and conversion parameters, moderately sensitive to subscription period and adoption rate, and strongly sensitive to audience distribution (a median of 500 followers eliminates it). The advantage persists for any CCRMS fee below approximately 17\%.

\subsection{Prototype-Bound versus Simulated Cohorts}\label{subsec:prototype-bound}

The win-rate results in Tables~7.9 and~7.10 characterize the CCRMS design. They do not characterize the artifact as benchmarked.

The prototype stores each first-time buyer as a child NFT within a \texttt{BoundedVec<\_, ConstU32<50>{}>}. Expired subscriptions do not free their slots. Consequently, a single content item can accumulate at most fifty subscribers and owners over its lifetime (Appendix~A.8). The simulation's Large, Huge, and Mega tiers assume no such cap. These are the tiers in which CCRMS wins (Table~7.9). Until the mitigation in Appendix~\ref{sec:e-child-cap} is implemented, the prototype cannot host the audiences for which the fee-side advantage is claimed.

What the current artifact supports is independent of that cap: the three rights maps and the priority \texttt{check\_access} (Claim~1, Section~\ref{subsec:claims}); the fee assumptions used in the simulation; and the issuance-style XCM on local HRMP (Claim~2). The economic figures remain a design-level result under Appendix~\ref{app:monte-carlo} (Claim~3, Section~\ref{subsec:claims}).

\section{Security and User-Centric Results}\label{sec:security-user-results}


\subsection{Security Audit Findings}\label{subsec:security-audit}

The security audit (Scenario 6, Section~\ref{sec:test-scenarios}) produced nine findings labeled A through I. Table~\ref{tab:security-findings} summarizes them.

\begin{longtable}{@{}L{0.04\linewidth} L{0.16\linewidth} L{0.50\linewidth} L{0.22\linewidth}@{}}
\caption{Security Audit Findings}\label{tab:security-findings}\\
\toprule
\textbf{ID} & \textbf{Severity} & \textbf{Description} & \textbf{Status} \\
\midrule
\endfirsthead

\multicolumn{4}{l}{\emph{Table~\ref{tab:security-findings} continued from previous page}}\\
\toprule
\textbf{ID} & \textbf{Severity} & \textbf{Description} & \textbf{Status} \\
\midrule
\endhead

\midrule
\multicolumn{4}{r}{\emph{Continued on next page}}\\
\endfoot

\bottomrule
\endlastfoot

A & Low & XCM extrinsics callable locally as gift transactions (intentional) & Accepted design \\
\addlinespace
B & \textbf{High} & Authorization gap in \texttt{xcm\_transfer\_ownership}: any funded account could transfer any user's ownership token & \textbf{Fixed} (commit \texttt{9ceca57}) \\
\addlinespace
C & Low & Permissive XCM filter configuration (\texttt{SafeCallFilter}, \texttt{XcmExecuteFilter}, \texttt{XcmTeleportFilter}, \texttt{XcmReserveTransferFilter} all \texttt{Everything}) & Accepted prototype config \\
\addlinespace
D & Low & Zero-price content registration allowed & Accepted feature \\
\addlinespace
E & Low & View pack overwrite on re-purchase & \textbf{Fixed} (commit \texttt{9ceca57}) \\
\addlinespace
F & Informational & Single-collator MTTR $\sim$15--20 seconds & Test environment artifact \\
\addlinespace
G & Positive & 100\% state persistence across crash & Confirmed \\
\addlinespace
H & Positive & Block production rate unchanged after recovery & Confirmed \\
\addlinespace
I & Medium (design limit) & Zero-view denial of service: \texttt{purchase\_views(num\_views=0)} charges \$0 but mints a child NFT, permanently exhausting a child slot & Open: unmitigated design-limit DoS; fix specified in Appendix~\ref{sec:e-child-cap} \\
\end{longtable}

Two findings (B and E) have been remediated; one (I) is an open Medium-severity design-limit denial-of-service that consumes a child slot at zero price (fix specified in Appendix~\ref{sec:e-child-cap}); three (A, C, D) are accepted design properties; and three (F, G, H) are positive reliability observations. \textbf{Finding~B} and \textbf{Finding~E} were addressed in commit \texttt{9ceca57}; the Finding~B remediation is conservative (the equality check cannot be satisfied by a remote XCM caller, so cross-chain transfers via \texttt{xcm\_transfer\_ownership} are pending a follow-on mitigation). Appendix~\ref{app:security} provides the complete narrative.

\subsection{Test Coverage}\label{subsec:test-coverage}

Our unit test suite covers \textbf{16 of the 19 distinct error variants} defined in \texttt{pallet-content-rights}, achieving 84 percent coverage, which exceeds the Phase~2 objective of 80 percent. Three variants remain untested and are out of scope for the unit test suite: \texttt{ItemIdOverflow} requires exhausting a \texttt{u32} counter under normal extrinsic dispatch, which is not feasible in a test environment; \texttt{NftOperationFailed} is triggered only by internal \texttt{pallet-nfts} state corruption that the mock runtime does not replicate; and \texttt{ChildAlreadyNested} cannot be reached through any publicly exposed extrinsic. The comprehensive error variant table is provided in Appendix~\ref{app:security}.

\subsection{Decentralization Assessment (HHI)}\label{subsec:decentralization-hhi}

The Herfindahl-Hirschman Index (HHI) was calculated from the Polkadot mainnet validator distribution (parachains inherit relay-chain security). Comparison values for Ethereum PoS and Solana are our estimates from publicly reported stake distributions (Rated.network; Solana Compass); the U.S.\ DOJ Horizontal Merger Guidelines define HHI below 1,500 as unconcentrated \citep{dojFTC2010}. Table~\ref{tab:hhi} presents the comparison.

\begin{longtable}{@{}L{0.28\linewidth} L{0.30\linewidth} R{0.15\linewidth} R{0.22\linewidth}@{}}
\caption{Herfindahl-Hirschman Index (HHI) Comparison}\label{tab:hhi}\\
\toprule
\textbf{System} & \textbf{Validators / Entities} & \textbf{HHI (approx.)} & \textbf{Nakamoto Coefficient} \\
\midrule
\endfirsthead

\multicolumn{4}{l}{\emph{Table~\ref{tab:hhi} continued from previous page}}\\
\toprule
\textbf{System} & \textbf{Validators / Entities} & \textbf{HHI (approx.)} & \textbf{Nakamoto Coefficient} \\
\midrule
\endhead

\midrule
\multicolumn{4}{r}{\emph{Continued on next page}}\\
\endfoot

\bottomrule
\endlastfoot

Polkadot mainnet (NPoS, $\sim$600 validators, as of September 2026) & $\sim$600 (near-equal stake) & $\sim$16.7 & $\sim$160--200 \\
\addlinespace
Ethereum PoS & $\sim$900K validators (concentrated in $\sim$5 staking providers) & $\sim$1,200--2,000 & $\sim$5--7 \\
\addlinespace
Solana & $\sim$1,900 validators (top-heavy stake distribution) & $\sim$200--500 & $\sim$19--25 \\
\addlinespace
Centralized DRM (e.g., Spotify) & 1 entity & 10,000 & 1 \\
\addlinespace
US DOJ ``unconcentrated'' threshold & -- & $<$1,500 & -- \\
\end{longtable}

Polkadot's HHI of $\sim$16.7 (as of September 2026, $\sim$600 validators) is well below the U.S.\ DOJ ``unconcentrated'' threshold of 1,500 and orders of magnitude lower than a centralized monopoly at 10,000. The Nakamoto coefficient ($\sim$160--200 for Polkadot vs.\ $\sim$5--7 for Ethereum PoS and 1 for centralized DRM) confirms that CCRMS on Polkadot exceeds the decentralization of all alternatives by approximately 600-fold.

Table~\ref{tab:hhi} should be read as a statement about safety under a production parachain, not about the test network. The relay-chain HHI indicates who secures CCRMS state once a block is included. The evaluated prototype's liveness is limited to a single collator (Nakamoto coefficient 1 for block production). A deployed CCRMS would inherit Polkadot's validator set for safety and would still need a collator set (or equivalent coretime) for liveness. The 600-fold HHI contrast with centralized DRM is therefore a property of the intended hosting environment, not of the Zombienet used in Chapter~\ref{chap:results}.

\subsection{Static Analysis and Dependency Audit}\label{subsec:static-analysis}

\texttt{cargo clippy} reported zero warnings. \texttt{cargo audit} reported 8 advisories in transitive SDK dependencies; none affect pallet code and all are mitigated by the sandboxed WASM runtime architecture (full list in Appendix~\ref{app:security}).

\section{Summary}\label{sec:chapter-7-summary}

Table~\ref{tab:kpi-results} consolidates the measured outcomes for all 11 KPIs.

\begin{longtable}{@{}L{0.03\linewidth} L{0.22\linewidth} L{0.18\linewidth} L{0.33\linewidth} L{0.18\linewidth}@{}}
\caption{KPI Results Summary}\label{tab:kpi-results}\\
\toprule
\textbf{\#} & \textbf{KPI} & \textbf{Target} & \textbf{Measured} & \textbf{Target met?} \\
\midrule
\endfirsthead

\multicolumn{5}{l}{\emph{Table~\ref{tab:kpi-results} continued from previous page}}\\
\toprule
\textbf{\#} & \textbf{KPI} & \textbf{Target} & \textbf{Measured} & \textbf{Target met?} \\
\midrule
\endhead

\midrule
\multicolumn{5}{r}{\emph{Continued on next page}}\\
\endfoot

\bottomrule
\endlastfoot

1 & Throughput & $>$10 TPS & \textbf{26.2 TPS} sustained & \textbf{Yes} \\
\addlinespace
2 & Inclusion latency & $<$12~s & \textbf{$\sim$6~s} (1 block) & \textbf{Yes} \\
\addlinespace
3 & XCM inclusion delay & $<$10 blocks & \textbf{0--7 blocks (median 5)} & \textbf{Yes} \\
\addlinespace
4 & Block utilisation & $<$75\% & \textbf{12.9\%} at 300 txs (reference time only) & \textbf{Yes --- utilisation headroom confirmed} \\
\addlinespace
5 & Storage efficiency & $<$500 bytes & \textbf{191 bytes}/content, \textbf{112 bytes}/subscriber & \textbf{Yes} \\
\addlinespace
6 & Uptime & $>$85\% & \textbf{90\%} & \textbf{Yes} \\
\addlinespace
7 & MTTR & $<$60~s & \textbf{$\sim$15--20~s} & \textbf{Yes} \\
\addlinespace
8 & State persistence & 100\% & \textbf{100\%} & \textbf{Yes} \\
\addlinespace
9 & Security findings & No Critical/High unremediated at submission & \textbf{1 High (Finding~B) identified and patched by the author; 0 unremediated Critical/High findings in that audit. No third-party review (Section~\ref{subsec:operational-limitations}). Finding~I remains an unmitigated design-limit DoS (Table~\ref{tab:security-findings}, Appendix~\ref{sec:e-child-cap}).} & Met as interpreted: no unremediated Critical/High findings in the author's audit. Not independently verified. \\
\addlinespace
10 & Test coverage & $>$80\% & \textbf{84\%} (16/19 error variants) & \textbf{Yes} \\
\addlinespace
11 & Decentralization ratio & Documented & \textbf{$270\times$ slower, qualitatively superior} & \textbf{Yes} \\
\end{longtable}

Ten of eleven KPIs met their Table~\ref{tab:kpi-defs} targets for the stated measurements. KPI~9 is recorded as met only under the interpretation ``no unremediated Critical/High in the author's audit''; there was no third-party review, and Finding~I remains open as a design-limit DoS (Section~\ref{sec:limitations}, Appendix~\ref{sec:e-child-cap}).
